%=============================================================================
% CROSS-REFERENCE: THE LATTICE beta -> kappa MAP IN DFD's U(1) THEORY
% Generated: 2026-03-22
% Source files: Ab_Initio_Derivation...pdf (Eqs. 11-13, 19-20),
%   appendix_K.tex, section_quantum.tex, section_alpha_locked.tex,
%   appendix_G.tex
%=============================================================================

\section{The $\beta \to \kappa$ Map: How DFD Extracts Physics from the Lattice}

\subsection{Executive Summary}

DFD does \textbf{not} use any standard analytic $\beta \to \kappa$ renormalization formula
(neither mean-field, nor perturbative, nor Toeplitz-based).
Instead, $\kappa$ is a \textbf{non-perturbative Monte Carlo measurement} extracted
from the lattice via the background-field method. The relationship
$\kappa(\beta)$ is determined entirely numerically; no closed-form expression
connects $\beta_{U(1)}$ to $\kappa_{U(1)}$.

\medskip
\noindent\textbf{Key distinction (PDF Eq.~11):}
\begin{equation}
\beta_{U(1)} = \frac{1}{g_1^2} \approx 3.80 \quad \text{(bare lattice input)},
\qquad
\kappa_{U(1)} \approx 7.25 \quad \text{(measured output)}.
\end{equation}
These are \emph{not} the same quantity. The paper explicitly states:
``The relationship between them is non-perturbative and lattice-size dependent,
which is precisely why the simulation is needed.''

%-----------------------------------------------------------------------------
\subsection{The Three Layers of the DFD Lattice Calculation}
%-----------------------------------------------------------------------------

The DFD $\alpha$-derivation involves three logically distinct layers:

\paragraph{Layer 1: Topological inputs $\to$ bare lattice couplings $\beta$.}
Four constraints, all derived from topology with zero continuous free parameters:
\begin{enumerate}
\item $k_{\max} = \chi(\mathbb{C}P^2, \mathcal{O}(9) \oplus \mathcal{O}^{\oplus 5}) = 60$
      \quad [Spin$^c$ index, Bridge Lemma]
\item $\beta_{U(1)} = \langle k_{\rm eff} \rangle_{k_{\max}=60}
      = \frac{\sum_{k=0}^{59}(k+2)\,w(k)}{\sum_{k=0}^{59}w(k)} = 3.7969 \approx 3.80$
      \quad [CS vacuum expectation]
\item $\beta_{SU(2)}/\beta_{U(1)} = (n_2/n_1) \times N_{\rm gen} = 2 \times 3 = 6$
      \quad [Frame Stiffness Theorem $\times$ index theorem]
\item $\beta_{SU(2)} = 6 \times 3.80 = 22.80$ \quad [derived]
\end{enumerate}
Here $w(k) = \frac{2}{k+2}\sin^2\!\bigl(\frac{\pi}{k+2}\bigr)$ are the SU(2) Chern--Simons
weights and $k_{\rm eff} \equiv k + 2$ is the shifted level (dual Coxeter number $h^\vee = 2$
for SU(2)).

\paragraph{Layer 2: Monte Carlo simulation $\to$ measured stiffnesses $\kappa$.}
The lattice action (PDF Eq.~19) couples the gauge sector to the microsector
via the $\psi(k)$ field:
\begin{equation}
S = \sum_x [-\log w(k_x)]
  + \frac{K_\psi}{2}\sum_{\langle xy \rangle}(\psi_x - \psi_y)^2
  - \sum_p \beta\, e^{-\psi_p} \cos(\theta_p + \theta_{\rm bg}\,\Omega_p),
\end{equation}
where $\psi_x = \psi(k_x)$ is the coarse-graining map and $\Omega_p$ marks
background-field plaquettes. Metropolis updates generate equilibrated
configurations at the input $(\beta_{U(1)}, \beta_{SU(2)})$.

The renormalized stiffness is then extracted via the \textbf{background-field method}
(PDF Eq.~20):
\begin{equation}\label{eq:kappa-extraction}
\boxed{\kappa = \frac{F''(0)}{V}}, \qquad
F''(0) = \langle S'' \rangle - \langle (S')^2 \rangle + \langle S' \rangle^2,
\end{equation}
where primes denote derivatives with respect to background field strength
$\theta_{\rm bg}$ evaluated at $\theta_{\rm bg} = 0$, and $V = L^3$.
This is the standard susceptibility/free-energy curvature method used in
lattice gauge theory. It captures all non-perturbative effects including
quantum fluctuations, confinement/deconfinement physics, and finite-size
corrections.

\paragraph{Layer 3: Measured stiffnesses $\to$ $\alpha$.}
Given the measured $\kappa_{U(1)}$ and $\kappa_{SU(2)}$, the gauge couplings are
extracted via Wilson's normalization conventions (PDF Eqs.~12--13):
\begin{align}
g_1^2 &= \frac{1}{\kappa_{U(1)}}, \qquad
g_2^2 = \frac{4}{\kappa_{SU(2)}}, \\[6pt]
\frac{1}{e^2} &= \frac{1}{g_1^2} + \frac{1}{g_2^2}
              = \kappa_{U(1)} + \frac{\kappa_{SU(2)}}{4}, \\[6pt]
\alpha_W &= \frac{e^2}{4\pi}
          = \frac{(1/\kappa_{U(1)})(4/\kappa_{SU(2)})}
                 {(1/\kappa_{U(1)}) + (4/\kappa_{SU(2)})}
            \cdot \frac{1}{4\pi}.
\end{align}
The factor of 4 in $g_2^2 = 4/\kappa_{SU(2)}$ is the standard Wilson action
normalization for SU(2): with $\mathrm{Tr}(T^a T^b) = \frac{1}{2}\delta^{ab}$,
the Wilson plaquette action gives $\beta_{SU(2)} = 4/g^2$.

%-----------------------------------------------------------------------------
\subsection{Why No Analytic $\beta \to \kappa$ Formula Exists}
%-----------------------------------------------------------------------------

In standard compact U(1) lattice gauge theory, the relationship between
the bare coupling $\beta = 1/g^2$ and the renormalized stiffness $\kappa$
is complicated by:

\begin{itemize}
\item \textbf{Confinement/Coulomb phase transition:} Compact U(1) on a
      finite lattice has a first-order phase transition near $\beta_c \approx 1.01$
      (4D). DFD operates at $\beta = 3.80$, well into the Coulomb phase,
      but finite-size effects from the lattice sizes $L = 6$--$16$ are
      non-negligible.

\item \textbf{Non-perturbative renormalization:} The perturbative expansion
      $\kappa = \beta(1 + c_1/\beta + \ldots)$ is an asymptotic series that
      does not converge at $\beta = 3.80$. Mean-field estimates
      $\kappa \approx \beta \langle P \rangle$ (where $\langle P \rangle$ is the
      plaquette expectation value) give only a rough approximation.

\item \textbf{DFD microsector coupling:} The DFD action (Eq.~19) is
      \emph{not} the standard compact U(1) Wilson action. The additional
      $\psi(k)$ field and its nearest-neighbor coupling $K_\psi$ modify
      the effective dynamics. No analytic result exists for this coupled system.
\end{itemize}

The numerical result $\kappa_{U(1)} \approx 7.25$ at $\beta_{U(1)} = 3.80$
represents a renormalization factor of $\kappa/\beta \approx 1.9$, which is
substantially larger than unity. This large ratio reflects the non-perturbative
quantum fluctuation enhancement in the Coulomb phase of the microsector-coupled
lattice theory.

%-----------------------------------------------------------------------------
\subsection{Role of the Toeplitz Truncation}
%-----------------------------------------------------------------------------

The Toeplitz truncation enters at \textbf{Layer 1}, not Layer 2. Specifically:

\begin{itemize}
\item The Toeplitz quantization at level $m = k + 3$ on $\mathbb{C}P^1$ sets
      $d = k + 4 = 64$ (for $k = k_{\max} = 60$).
\item This determines the BOOST factor
      $\varepsilon_{\rm adj} = d^2/(d^2 - 1) = 4096/4095$ in the
      convention-locked $\alpha$ derivation (Section~8C of the unified review).
\item It does \emph{not} enter the $\beta \to \kappa$ map.
      The lattice simulation at Layer 2 knows nothing about the Toeplitz
      truncation; it simply runs at the input $\beta$ values and measures
      $\kappa$ from fluctuations.
\end{itemize}

The spectral cutoff $\Lambda^3 = k \cdot (k+3)/(k+4)$ from the Toeplitz
truncation enters only in the analytical (non-lattice) derivation path
that achieves sub-ppm precision (Section~8C, $\alpha^{-1} = 137.03599985$).
The lattice path achieves $\sim 1\%$ precision independently.

%-----------------------------------------------------------------------------
\subsection{Comparison with Standard Lattice Gauge Theory}
%-----------------------------------------------------------------------------

\begin{center}
\renewcommand{\arraystretch}{1.3}
\begin{tabular}{lcc}
\toprule
\textbf{Feature} & \textbf{Standard lattice U(1)} & \textbf{DFD lattice} \\
\midrule
Action & Wilson plaquette & Wilson + microsector $\psi(k)$ \\
$\beta$ meaning & $1/g^2_{\rm bare}$ & $\langle k_{\rm eff} \rangle$ from CS vacuum \\
$\beta$ value & Free parameter & Fixed: $3.80$ from topology \\
$\kappa$ extraction & Background field / susceptibility & Same method (Eq.~20) \\
$\beta \to \kappa$ map & Non-perturbative & Non-perturbative \\
Analytic approx. & $\kappa \approx \beta \langle P \rangle$ & None available \\
Measured $\kappa/\beta$ & $\sim 1$ (large $\beta$) & $\approx 1.9$ \\
Physical observable & $g_{\rm renorm}^2 = 1/\kappa$ & $\alpha_W$ via electroweak mixing \\
\bottomrule
\end{tabular}
\end{center}

%-----------------------------------------------------------------------------
\subsection{Answer to the Three Posed Questions}
%-----------------------------------------------------------------------------

\begin{enumerate}
\item \textbf{Is it $\kappa = \beta \times \langle P \rangle$ (mean-field)?}

No. DFD does not use a mean-field approximation. The stiffness $\kappa$
is measured non-perturbatively from the free-energy curvature (Eq.~20).
The measured ratio $\kappa_{U(1)}/\beta_{U(1)} \approx 7.25/3.80 \approx 1.9$
is much larger than typical mean-field plaquette values.

\item \textbf{Is it $\kappa = \beta \times (1 + c_1/\beta + \ldots)$ (perturbative)?}

No. No perturbative expansion is used. The paper explicitly states that
the $\beta \to \kappa$ relationship is ``non-perturbative and lattice-size
dependent, which is precisely why the simulation is needed.''

\item \textbf{Does it involve the Toeplitz truncation?}

No, not at the $\beta \to \kappa$ level. The Toeplitz truncation determines
the input $\beta$ (Layer 1) and the BOOST factor in the analytical path,
but does not enter the lattice measurement of $\kappa$ (Layer 2).
\end{enumerate}

%-----------------------------------------------------------------------------
\subsection{Summary}
%-----------------------------------------------------------------------------

\begin{tcolorbox}[colback=blue!5!white, colframe=blue!50!black,
  title=The $\beta \to \kappa$ Map in DFD]

\textbf{Method:} Non-perturbative Monte Carlo measurement via
background-field free-energy curvature.

\textbf{Formula:}
$\kappa = F''(0)/V$ where $F''(0) = \langle S'' \rangle
- \langle (S')^2 \rangle + \langle S' \rangle^2$

\textbf{Key numbers:}
\begin{itemize}[nosep]
\item $\beta_{U(1)} = 3.80$ (bare input from CS vacuum)
\item $\kappa_{U(1)} \approx 7.25$ (measured output)
\item $\kappa_{U(1)}/\kappa_{SU(2)} \approx 0.50$ (confirms Frame Stiffness Theorem)
\item $\alpha_W \approx 1/137$ (extracted via electroweak mixing)
\end{itemize}

\textbf{Not used:} Mean-field, perturbative expansion, or Toeplitz-based
analytic continuation. The entire $\beta \to \kappa$ map is numerical.

\textbf{Why this matters:} The simulation is the irreducible non-perturbative
step that converts topological inputs ($\beta$) into physical outputs ($\alpha$).
No analytic shortcut exists for the DFD microsector-coupled action.
\end{tcolorbox}

\end{document}
